% !TeX root = ./main.tex

% --------------------------------------------------
% 資訊設定（Information Configs）
% --------------------------------------------------

\ntusetup{
  university*   = {National Taiwan University},
  university    = {國立臺灣大學},
  college       = {電機資訊學院},
  college*      = {Electrical Engineering and Computer Science},
  institute     = {電信工程學研究所},
  institute*    = {Graduate Institute of Communication Engineering College},
  title         = {上行非正交多重存取接收器之效能分析與基於調變編碼方式之使用者配對應用},
  title*        = {Performance Analysis of Uplink NOMA Receivers and Their Application to MCS-Based User Pairing},
  author        = {洪睿帆},
  author*       = {Jui-Fan Hung},
  ID            = {R13942124},
  advisor       = {林茂昭, 李世凱},
  advisor*      = {Mao-Chao Lin, Shih-Kai Lee},
  date          = {2026-07-28},         % 若註解掉，則預設為當天
  oral-date     = {2026-07-28},         % 若註解掉，則預設為當天
  DOI           = {10.5566/NTU2018XXXXX},
  keywords      = {非正交的多重存取接收器, 多資源功率域非正交的多重存取, 分增益多重存取, 調變編碼方式之使用者配對}, %關鍵字中文
  keywords*     = {NOMA rceiver, multi-resource PD-NOMA, GDMA, MCS user pairing}, %關鍵字英文
}

% --------------------------------------------------
% 加載套件（Include Packages）
% --------------------------------------------------

\usepackage[numbers,sort&compress]{natbib}      % 參考文獻
\usepackage{amsmath, amsthm, amssymb}   % 數學環境
\usepackage{amssymb, amsfonts, bm}

\usepackage{subcaption}
\usepackage{textcomp}
\usepackage{xcolor}
\usepackage{tikz}
\usepackage{physics} % 引入 physics 套件

\usepackage{algorithm}
\usepackage{algorithmic}
%\usepackage{algorithmicx}
%\usepackage{algpseudocode}

\usepackage{ulem, CJKulem}              % 下劃線、雙下劃線與波浪紋效果
\usepackage{booktabs}                   % 改善表格設置
\usepackage{multirow}                   % 合併儲存格
\usepackage{diagbox}                    % 插入表格反斜線
\usepackage{array}                      % 調整表格高度
\usepackage{longtable}                  % 支援跨頁長表格
\usepackage{paralist}                   % 列表環境


\usetikzlibrary{arrows.meta, shapes.geometric, calc}
\usepackage{booktabs}
\usepackage[table]{xcolor}
\usepackage{nicematrix}
\usepackage{makecell}

\usepackage{lipsum}                     % 英文亂字
\usepackage{zhlipsum}                   % 中文亂字
\usepackage{comment}
\usepackage{graphicx}
% --------------------------------------------------
% 套件設定（Packages Settings）
% --------------------------------------------------
\usetikzlibrary{positioning, arrows.meta, fit, calc, decorations.pathreplacing}
